\subsection{Human Motion and Gait Analysis}

Human gait has long been studied as a compact yet informative representation of mobility, functional status, and individual motion characteristics.
Early computer-vision work treated walking as a spatio-temporal signal rather than a collection of isolated poses.
For example, Ohara et al. introduced the concept of \emph{gait volume}, in which image sequences of walking are represented as a spatio-temporal volume and analyzed in the frequency domain to capture person-specific walking rhythms~\cite{ohara2004gaitvolume}.
This line of work is relevant to our study because it highlights that aging-related motion signatures are not only spatial, e.g., posture or limb configuration, but also temporal, e.g., rhythm, cadence, and phase relationships across body segments.

In parallel, clinical gait analysis has provided a rigorous foundation for quantifying human mobility through interpretable biomechanical and spatio-temporal metrics.
Parameters such as walking speed, step length, stride length, cadence, double-support time, and gait variability are commonly used to characterize mobility and functional decline across adulthood~\cite{herssens2020investigation}.
Recent sensor-based and markerless approaches further suggest that gait analysis can be moved beyond specialized motion-capture laboratories, using depth cameras or wearable sensors to estimate clinically meaningful gait parameters in more accessible settings~\cite{wagner2023methods}.
These developments are important for assistive robotics, where robots operating in homes, clinics, or care facilities must reason about human mobility from practical sensing setups rather than ideal laboratory recordings.

Learning-based skeleton representations have further expanded the scope of motion analysis.
Spatial-Temporal Graph Convolutional Networks (ST-GCNs) model the human body as a graph whose nodes correspond to joints and whose edges encode anatomical and temporal dependencies~\cite{yan2018spatial}.
Such graph-based models are well suited for analyzing gait and full-body movement because they preserve both the structured connectivity of the body and the temporal evolution of motion.
However, most gait-analysis and skeleton-recognition studies remain primarily discriminative: they classify actions, estimate identity, detect abnormality, or predict demographic attributes.
They do not directly provide generative models that can synthesize controllable, age-conditioned 3D motion for downstream robotics applications.
Our work addresses this gap by using clinical gait recordings not only for recognizing age-related motion signatures, but also for guiding age-aware human motion generation.

\subsection{Motion Generation for Human-Centered Robotics}

Recent progress in human motion generation has been driven by large-scale motion-capture archives and text-conditioned generative models.
Datasets such as AMASS have unified heterogeneous motion-capture collections into parametric body representations~\cite{mahmood2019amass}, while text-motion datasets such as HumanML3D have enabled natural-language-conditioned 3D motion synthesis~\cite{guo2022generating}.
Building on these resources, diffusion-based models have become a major paradigm for motion generation.
The Motion Diffusion Model (MDM) adapts classifier-free diffusion to human motion and supports text-to-motion, action-to-motion, and motion editing tasks~\cite{tevet2023human}.
Motion Latent Diffusion (MLD) improves efficiency by performing diffusion in a learned motion latent space~\cite{chen2023executing}, while MotionDiffuse demonstrates fine-grained text-driven generation through iterative denoising~\cite{zhang2024motiondiffuse}.
These methods show that deep generative models can synthesize diverse and semantically meaningful human motions.

For human-centered robotics, however, the goal is not only to generate visually plausible motion, but also to model the kinds of human movement that robots must anticipate, assist, or safely accommodate.
Human-aware planning research has shown that robots can behave more safely and efficiently when they explicitly account for human motion, personal space, and future human trajectories~\cite{sisbot2007human,unhelkar2018human}.
Assistive robotics further strengthens this need because the human partner may have age-related or condition-related mobility constraints.
Simulation frameworks such as Assistive Gym explicitly model human motion, physical capabilities, and preferences in activities of daily living~\cite{erickson2020assistivegym}.
Similarly, the HARMONIC dataset provides multimodal observations of human--robot interaction during assistive eating, supporting research on intention prediction, shared autonomy, and human-state modeling~\cite{newman2022harmonic}.

Despite this progress, most existing human motion generation models are trained on general-purpose motion datasets dominated by young or healthy actors, and their conditioning mechanisms usually focus on action labels, language prompts, trajectories, or style.
Population-level biomechanical factors, such as age, functional mobility, and clinical gait characteristics, are rarely modeled explicitly.
This limitation is particularly important for robots intended to support aging societies, where safe and proactive assistance requires motion priors that reflect how older adults actually move.
Our work contributes toward this direction by integrating clinical gait analysis, age-group recognition, and text-conditioned motion diffusion to study whether generated motions can preserve measurable aging-related kinematic signatures.
